% docs/physics/latex/04_pion.tex - SIMC++ physics note 04: exclusive pion electroproduction.
% Standalone PDF docs/physics/pdf/SIMCpp_Physics_04_pion.pdf (copied to src/reactions/pion/ and src/physics/pion/); a
% chapter of the book SIMC++ Physics Reference. Rules: docs/physics/README.md; structure: template_reaction.tex.
\documentclass{simcnote}
\simccommon
\begin{document}
\simcnotemeta{04}{Exclusive (e,e\pr \texorpdfstring{$\pi$}{pi})}{Exclusive pion electroproduction: H, D, A(e,e\pr \texorpdfstring{$\pi^\pm$}{pi+-}) and (e,e\pr \texorpdfstring{$\pi^0$}{pi0})}%
  {Reaction plug-in \texttt{pion\_production} with the physics model \texttt{pion}: quasi-free and coherent
   charged and neutral pion production, nucleon and \texorpdfstring{$\Delta(1232)$}{Delta(1232)} recoil, Fermi motion, \texorpdfstring{$\pi^0\to\gamma\gamma$}{pi0 decay} and pion decay in flight}
\simcnotemaketitle
\simcabstract{Exclusive electroproduction of a single pion, detected in coincidence with the scattered electron, on
hydrogen, deuterium and heavier nuclei. The electron energy and both pairs of slopes are generated; for D and $A\ge3$ the
struck nucleon also gets a Fermi momentum and a missing energy, and the pion energy follows from two-body kinematics
with a nucleon, a $\Delta(1232)$ of Breit--Wigner mass or the final nucleus as recoil. The cross section is the 2021 fit
of P. Bosted to exclusive $\pi^\pm$ data, joined to the MAID2007 tables below $W=2$ GeV, transformed from
$\mathrm{d}\sigma/\mathrm{d}t\,\mathrm{d}\varphi^*$ to the laboratory with an exact jacobian and multiplied by the
virtual-photon flux. A $\pi^0$ is decayed into two photons for a calorimeter, a charged pion can decay in flight in the
spectrometer. Notation, glossary and common physics: note 00.}
\simcnotetoc

% =====================================================================================================================
\section{The reaction in SIMC++}\lbl{sec:reaction}

\subsection{Process and scope}
The reaction is $e + N \to e' + \pi + X$ on a nucleon $N$ (free for hydrogen, bound in D or a heavier nucleus), with the
pion detected in the hadron arm. The recoil $X$ is not detected. \tabr{tab:finalstates} lists the final states: one
input integer, \code{which\_pion}, chooses the charge of the pion, the struck nucleon and the recoil (nucleon,
$\Delta(1232)$, or the whole nucleus for coherent production), and \code{doing\_pizero} turns the charged pion into a
$\pi^0$ (mass $M_{\pi^0}=134.9766$ MeV instead of $M_\pi=139.57018$ MeV).

\begin{simctable}{footnotesize}{L{0.08\linewidth}L{0.07\linewidth}L{0.31\linewidth}L{0.09\linewidth}L{0.17\linewidth}L{0.28\linewidth}}
\caption{Final states (input \code{which\_pion}, $\pi^0$ = \code{doing\_pizero}): struck mass $M_N$, recoil mass $M_X$ and the number of
target nucleons $n_\mathrm{tgt}$ of the weight (one final state per row).}\lbl{tab:finalstates}\\
\toprule\rowcolor{simchead}\simcth{value} & \simcth{$\pi^0$} & \simcth{process} & \simcth{$M_N$} & \simcth{$M_X$} & \simcth{targets, $n_\mathrm{tgt}$}\\\midrule\endfirsthead
\toprule\rowcolor{simchead}\simcth{value} & \simcth{$\pi^0$} & \simcth{process} & \simcth{$M_N$} & \simcth{$M_X$} & \simcth{targets, $n_\mathrm{tgt}$}\\\midrule\endhead
\bottomrule\endlastfoot
0 & no & $\gamma^*p\to\pi^+n$ & $M_p$ & $M_n$ & H, D, A; $Z$\\
1 & no & $\gamma^*n\to\pi^-p$ & $M_n$ & $M_p$ & D, A (refused on H); $N$\\
2 & no & $\gamma^*p\to\pi^+\Delta^0$ (D: also $\gamma^*n\to\pi^+\Delta^-$) & $M_p$ & $M_\Delta$ (sampled) & H, D, A; \eqr{eq:cgfactors}\\
3 & no & $\gamma^*p\to\pi^-\Delta^{++}$ (D: also $\gamma^*n\to\pi^-\Delta^+$) & $M_p$ & $M_\Delta$ (sampled) & H, D, A; \eqr{eq:cgfactors}\\
10 & no & $\gamma^*A\to\pi^+A'$ (coherent) & $M_A$ & $M_A-M_p+M_n$ & $A\ge3$ ($A\neq4$); $Z$\\
11 & no & $\gamma^*A\to\pi^-A'$ (coherent) & $M_A$ & $M_A-M_n+M_p$ & $A\ge5$; $N$\\
0 & yes & $\gamma^*p\to\pi^0p$ & $M_p$ & $M_p$ & H, D, A; $Z$\\
1 & yes & $\gamma^*n\to\pi^0n$ & $M_n$ & $M_n$ & D, A (refused on H in the default mode); $N$\\
2 & yes & $\gamma^*p\to\pi^0\Delta^+$ & $M_p$ & $M_\Delta$ (sampled) & H, D, A; \eqr{eq:cgfactors}\\
3 & yes & $\gamma^*n\to\pi^0\Delta^0$ & $M_n$ & $M_\Delta$ (sampled) & D, A (refused on H); \eqr{eq:cgfactors}\\
\end{simctable}

What is simulated: the quasi-free production on a nucleon of momentum $\vec p_F$ (D and $A\ge3$) or at rest
(H, coherent production), the unpolarised cross section of \secr{sec:xsec} in the photon--nucleon centre-of-mass frame,
the $\Delta(1232)$ recoil mass distribution, the event chain of note 00 (radiation on the three legs, energy loss and
multiple scattering of the pion, Coulomb correction of the electron), the decay $\pi^0\to\gamma\gamma$ and the decay
in flight of a charged pion inside the spectrometer. What is not simulated: final-state interactions of the pion or
the recoil (no pion absorption or transparency factor in the nucleus, no Pauli blocking), a cross-section model for
nuclei other than the nucleon fit (coherent production uses the nucleon model with $M_N=M_A$), the decay of the
$\Delta$, and polarisation observables (the MAID polarised terms are computed and discarded).

\subsection{Components}
\begin{simctable}{footnotesize}{L{0.15\linewidth}L{0.16\linewidth}L{0.15\linewidth}L{0.25\linewidth}L{0.29\linewidth}}
\caption{The SIMC++ components of exclusive pion production (one component per row).}\lbl{tab:components}\\
\toprule\rowcolor{simchead}\simcth{role} & \simcth{interface} & \simcth{registered name} & \simcth{class} & \simcth{file}\\\midrule\endfirsthead
\toprule\rowcolor{simchead}\simcth{role} & \simcth{interface} & \simcth{registered name} & \simcth{class} & \simcth{file}\\\midrule\endhead
\bottomrule\endlastfoot
reaction & \code{IReaction} & \code{pion\_production} & \code{reactions::PionProduction} & \file{reactions/pion/PionProduction.cpp}\\
recoil-mass table & (class) & -- & \code{reactions::DeltaMassTable} & \file{reactions/pion/PionProduction.cpp}\\
event completion, recon & (base class) & -- & \code{reactions::EepReactionBase} & \file{reactions/EepReactionBase.cpp}\\
centre-of-mass transform & (function) & -- & \code{reactions::transformToCm} & \file{reactions/EventChain.cpp}\\
cross section & \code{IPhysicsModel} & \code{pion} & \code{pion::PionModel} & \file{physics/pion/PionModel.cpp}\\
cross-section kernels & (functions) & -- & \code{pion::peepi}, \code{sigParam2021}, \code{exclfit} & \file{physics/pion/PionCrossSections.cpp}\\
MAID tables & (class, function) & -- & \code{pion::MaidTables}, \code{pion::sigmaid} & \file{physics/pion/PionCrossSections.cpp}\\
unused fits & (functions) & -- & \code{pion::sigBlok}, \code{sigParam04}, \code{sigParam3000} & \file{physics/pion/PionCrossSections.cpp}\\
unused multipole model & (class) & -- & \code{pion::SigMultipole} & \file{physics/pion/SigMultipole.cpp}\\
$\pi^0$ decay & (function) & -- & \code{pion::pizeroDecay} & \file{physics/pion/PizeroDecay.cpp}\\
$\pi^0$ photons in a calorimeter & (class) & -- & \code{driver::PizeroCalorimeter} & \file{driver/PizeroCalorimeter.cpp}\\
missing energy, $A\ge3$ & (class) & -- & \code{nuclear::SfTable} & \file{physics/nuclear/SpectralFunctionTables.cpp}\\
decay in flight & (functions) & -- & \code{optics::project}, \code{CosyEngine::transp} & \file{optics/CosyEngine.cpp}\\
beam, target, radiation, arms, outputs & as in note 00 & \code{standard}, \code{hallc\_cryo}, \code{simc\_rad}, \ldots & & \ctabr{tab:plugins}\\
\end{simctable}

\subsection{Selection and override}
The input \code{doing\_pion} = 1 selects \code{pion\_production} with the model \code{pion} (\ctabr{tab:defaults});
\code{doing\_pion} together with \code{doing\_semi} selects semi-inclusive production instead. The set-up derives the
case from the target: H for $\mathrm{nint}(A)=1$, D for $\mathrm{nint}(A)=2$, $A\ge3$ otherwise, and coherent production
(treated like hydrogen, with $M_N=M_A$) for \code{which\_pion} $\ge10$. It sets $M_h$, $M_N$ and $M_X$ of
\tabr{tab:finalstates}, forces $M_A=M_p$ for H and $M_A=M_d$, $M_\mathrm{rec}=M_p+M_n-M_N$ for D, and replaces the
input recoil mass of coherent production by $M_A+M_N^\mathrm{struck}-M_N^\mathrm{recoil}$ when the two differ by more than
100 MeV. Inputs that the Fortran stops on (π$^-$ on hydrogen, coherent production on H, D, He-4 and coherent π$^-$ on
He-3, $\gamma^*n\to\pi^0\Delta^0$ on hydrogen, \code{pizero\_ngamma} other than 1 or 2) stop SIMC++ too; the
corrected mode also refuses π$^0$ from the neutron on hydrogen. The missing-mass cut window is opened to $\pm10^6$ MeV.

\begin{simctable}{footnotesize}{L{0.3\linewidth}L{0.38\linewidth}L{0.32\linewidth}}
\caption{JSON keys of this reaction and its model (one key per row).}\lbl{tab:json}\\
\toprule\rowcolor{simchead}\simcth{key} & \simcth{effect} & \simcth{example}\\\midrule\endfirsthead
\toprule\rowcolor{simchead}\simcth{key} & \simcth{effect} & \simcth{example}\\\midrule\endhead
\bottomrule\endlastfoot
\code{pipeline.reaction.data\_dir} & directory of \code{delta\_relativistic\_bw.inp} (default \code{data/pion}) & \code{"data/pion"}\\
\code{pipeline.reaction.verbose} & print the Fortran warnings of the reaction & \code{true}\\
\code{pipeline.physics.data\_dir} & directory of the MAID tables (default \code{data/pion}) & \code{"data/pion"}\\
\code{pipeline.physics.verbose} & print the warnings of the transform & \code{true}\\
\end{simctable}
A replacement model gets the completed vertex event and the main record (Fermi momentum, $\varepsilon$, $-t$,
$\varphi_{pq}$) and must return $\mathrm{d}^5\sigma/\mathrm{d}E'\mathrm{d}\Omega_e\mathrm{d}\Omega_\pi$ per nucleon in
$\mu$b/(MeV sr$^2$) (\secr{sec:require}), for example
\code{\{"pipeline": \{"reaction": "pion\_production", "physics": \{"type": "pion", "data\_dir": "/my/maid"\}\}\}}.

\subsection{Event flow}
\begin{simctable}{footnotesize}{L{0.05\linewidth}L{0.33\linewidth}L{0.62\linewidth}}
\caption{One try of exclusive pion production (one step per row; the general chain is \ctabr{tab:chain}).}\lbl{tab:flow}\\
\toprule\rowcolor{simchead}\simcth{\#} & \simcth{C++ call} & \simcth{what is specific to this reaction}\\\midrule\endfirsthead
\toprule\rowcolor{simchead}\simcth{\#} & \simcth{C++ call} & \simcth{what is specific to this reaction}\\\midrule\endhead
\bottomrule\endlastfoot
1 & \code{StandardBeam::generate} & nothing\\
2 & \code{PionProduction::generate} & both slope pairs, $E'$, and for D/$A\ge3$ the Fermi momentum and $E_m$ (\secr{sec:gen})\\
3 & \code{PionProduction::completeEv} & $M_\Delta$ for \code{which\_pion} 2/3, the pion energy (\eqr{eq:quad}), $W$, $\varepsilon$, $-t$, $\varphi_{pq}$, the $\pi^0$ photons\\
4 & \code{SimcRadiator::radiate} & tail-1 limit from the electron-energy range (\eqr{eq:tail1}); a tail-1 photon calls \code{completeEv} again\\
5 & \code{MonteCarlo::run} & the pion crosses the hadron arm and may decay (\secr{sec:decay}); for a $\pi^0$ the two photons go to the calorimeter\\
6 & \code{EepReactionBase::completeRecon} & $E_m^\mathrm{rec}$, missing masses and $-t^\mathrm{rec}$ (\secr{sec:recon})\\
7 & \code{PionProduction::completeMain} & $S=1$; \code{PionModel::sigccMain}, the factors of \eqr{eq:cgfactors}, $n_\mathrm{tgt}$, the weight\\
\end{simctable}

% =====================================================================================================================
\section{Kinematics, generated variables and generation volume}\lbl{sec:gen}
\subsection{Generated variables}
\begin{simctable}{small}{L{0.38\linewidth}L{0.31\linewidth}L{0.31\linewidth}}
\caption{Generated variables (one variable per row).}\lbl{tab:generated}\\
\toprule\rowcolor{simchead}\simcth{variable} & \simcth{H, coherent} & \simcth{D, $A\ge3$}\\\midrule\endfirsthead
\bottomrule\endlastfoot
electron slopes $y'_e$, $x'_e$ & flat & flat\\
pion slopes $y'_h$, $x'_h$ & flat & flat\\
electron energy $E'$ & flat & flat\\
Fermi momentum $p_F$, direction $\hat p_F$ & 0 & from the momentum distribution, isotropic\\
missing energy $E_m$ & 0 & D: binding energy; $A\ge3$: from the spectral function\\
recoil mass $M_X$ & fixed, or sampled for $\Delta$ & fixed, or sampled for $\Delta$\\
pion energy $E_h$ & solved (\eqr{eq:quad}) & solved (\eqr{eq:quad})\\
\end{simctable}
The four slopes are drawn flat in their generation ranges, which are the edges of note 00, and converted to angles with
\ceqr{eq:physang}. The electron energy is drawn flat inside the generated-energy range narrowed by the summed-energy limits
$\Sigma$ (for this reaction $\Sigma$ limits the electron energy alone):
\begin{feqa}{reactions/pion/PionProduction.cpp: PionProduction::generate}{event.f:276-286, 300-315}
x' &= x'_\mathrm{min}+\xi\,(x'_\mathrm{max}-x'_\mathrm{min})\quad\text{(each of }y'_e, x'_e, y'_h, x'_h\text{)},\nonumber\\
E'^\mathrm{lo} &= \max\bigl(E'^\mathrm{gen}_\mathrm{min},\Sigma_\mathrm{min}\bigr),\qquad
E'^\mathrm{hi} = \min\bigl(E'^\mathrm{gen}_\mathrm{max},\Sigma_\mathrm{max}\bigr),\nonumber\\
E' &= E'^\mathrm{lo}+\xi\,(E'^\mathrm{hi}-E'^\mathrm{lo}),\qquad
w_\mathrm{gen}\leftarrow w_\mathrm{gen}\,\frac{E'^\mathrm{hi}-E'^\mathrm{lo}}{E'^\mathrm{gen}_\mathrm{max}-E'^\mathrm{gen}_\mathrm{min}},
\lbl{eq:genee}
\end{feqa}
the try failing when $E'^\mathrm{lo}>E'^\mathrm{hi}$. The pion energy is not generated.

\subsection{Fermi momentum and missing energy (D and \texorpdfstring{$A\ge3$}{A>=3})}
The momentum distribution is a table of nodes $p_i$ and cumulative probabilities $c_i$ ($i=1\ldots n$, normalised to
$c_n=1$ when it is read): \code{deut.dat} for D, \code{he3.dat}, \code{he4.dat}, \code{c12.dat} for $A=3$, 4, 12, and
\code{c12.dat} for every other $A$. With five uniform numbers:
\begin{feqa}{reactions/pion/PionProduction.cpp: PionProduction::generate}{event.f:339-360; dbase.f:582-589}
i &= \min\{i:\ \xi_1\le c_i\}\ (\text{or } n),\qquad
p_\mathrm{lo} = \tfrac12(p_{i-1}+p_i)\ (0\text{ for } i=1),\nonumber\\
p_\mathrm{hi} &= \tfrac12(p_i+p_{i+1})\ (p_n\text{ for } i=n),\qquad
p_F = p_\mathrm{lo}+\xi_2\,(p_\mathrm{hi}-p_\mathrm{lo}),\nonumber\\
\hat p_F &= (\sin\theta_F\cos\varphi_F,\ \sin\theta_F\sin\varphi_F,\ \cos\theta_F),\qquad
\cos\theta_F = 2\xi_3-1,\quad \varphi_F = 2\pi\xi_4,
\lbl{eq:fermi}
\end{feqa}
in the laboratory frame of note 00. The missing energy and the energy of the struck nucleon are
\begin{feqa}{reactions/pion/PionProduction.cpp: PionProduction::generate}{event.f:362-372}
\text{D:}\quad E_m &= M_p+M_n-M_d,\qquad \text{$A\ge3$:}\quad E_m \sim S(E_m, p_F)\ \text{at fixed } p_F,\nonumber\\
M_\mathrm{spec} &= M_A-M_N+E_m,\qquad E_F = M_A-\sqrt{M_\mathrm{spec}^2+p_F^2},
\lbl{eq:efer}
\end{feqa}
and $E_F=M_N$, $p_F=0$, $E_m=0$ for H and coherent production. For $A\ge3$ the missing energy is drawn from the Benhar
spectral-function table of note 02 (\code{benharsf\_3mod.dat}, \code{\_4}, \code{\_12}, \code{\_56}, \code{\_64},
\code{\_197}, otherwise \code{benharsf\_12.dat}; the proton column for $M_N=M_p$, the neutron column for $M_N=M_n$):
\begin{feq}{physics/nuclear/SpectralFunctionTables.cpp: SfTable::generateEm}{sf\_lookup.f:216-237}
C_j = \sum_{k\le j} S(E_k,p_F)\Big/\sum_{k} S(E_k,p_F),\qquad
j = \min\{j:\ \xi_5\le C_j\},\qquad E_m = E_j + \Delta E_j\,(\xi_6-\tfrac12),
\lbl{eq:genem}
\end{feq}
with $S(E,p)$ interpolated in the table as in note 02 and $E_j$, $\Delta E_j$ the centres and widths of the energy
bins. The momentum and the missing energy come from two different tables (the momentum distribution and the spectral
function), so $p_F$ and $E_m$ are not drawn from one joint distribution. The energy $E_F$ of \eqr{eq:efer} is the
nucleon's share of the target mass with the spectator on its mass shell, so the struck nucleon is off shell.

\subsection{Limits}
The vertex limits, the largest photon energy and the summed-energy limits are, with $M_X$ the recoil mass at set-up
time (the nominal $M_\Delta=1232$ MeV for \code{which\_pion} 2/3):
\begin{feqa}{setup/SetupLimits.cpp: BuildLimits}{init.f:344-358, 380-402}
\text{H, coherent:}&\quad E_m = P_m = 0,\qquad M_\mathrm{rec}=T_\mathrm{rec}=0,\nonumber\\
\text{D:}&\quad E_m = M_p+M_n-M_d,\qquad P_m\in[0,\ p_n],\nonumber\\
A\ge3\text{:}&\quad E_m\in[M_N+M_\mathrm{rec}-M_A,\ E_{m}^\mathrm{table,max}],\qquad P_m\in[0,\ p_n],\nonumber\\
\text{D, }A\ge3\text{:}&\quad M_\mathrm{rec}^\mathrm{min,max}=M_A-M_N+E_m^\mathrm{min,max},\nonumber\\
&\quad T_\mathrm{rec}^\mathrm{min}=\sqrt{(M_\mathrm{rec}^\mathrm{max})^2+(P_m^\mathrm{min})^2}-M_\mathrm{rec}^\mathrm{max},\quad
T_\mathrm{rec}^\mathrm{max}=\sqrt{(M_\mathrm{rec}^\mathrm{min})^2+(P_m^\mathrm{max})^2}-M_\mathrm{rec}^\mathrm{min},\nonumber\\
T_X^\mathrm{max} &= E_b^\mathrm{max}+M_N-M_X-E'_\mathrm{min}-E_{h,\mathrm{min}}-E_m^\mathrm{min}-T_\mathrm{rec}^\mathrm{min},\qquad T_X^\mathrm{min}=0,
\lbl{eq:vertexedge}
\end{feqa}
($p_n$ the last node of the momentum distribution, $T_X$ the kinetic energy of the recoil), and
\begin{feqa}{setup/SetupLimits.cpp: BuildLimits}{init.f:411-425, 456-469, 477-488, 500-513}
E_\gamma^\mathrm{tot,max} &= E_b^\mathrm{max}+M_N-M_X-E'_\mathrm{min}-E_{h,\mathrm{min}}-E_m^\mathrm{min}-T_\mathrm{rec}^\mathrm{min},\nonumber\\
\Sigma_\mathrm{max} &= \min\bigl(E_b^\mathrm{max}+M_N-M_X-E_{h,\mathrm{min}}-E_m^\mathrm{min}-T_\mathrm{rec}^\mathrm{min},\ E'_\mathrm{max}+E_{\gamma2}^\mathrm{max}\bigr)+d_E,\nonumber\\
\Sigma_\mathrm{min} &= \max\bigl(0,\ \max\bigl(E_b^\mathrm{min}+M_N-M_X-E_{h,\mathrm{max}}-E_m^\mathrm{max}-T_\mathrm{rec}^\mathrm{max}-T_X^\mathrm{max}-E_\gamma^\mathrm{tot,max},\ E'_\mathrm{min}\bigr)-d_E\bigr),\nonumber\\
E'^\mathrm{gen} &= \bigl[\max(\Sigma_\mathrm{min},E'_\mathrm{min}),\ \min(\Sigma_\mathrm{max},E'_\mathrm{max}+E_{\gamma2}^\mathrm{max})\bigr],\qquad
E_h^\mathrm{gen} = \bigl[E_{h,\mathrm{min}},\ E_{h,\mathrm{max}}+E_{\gamma3}^\mathrm{max}\bigr],
\lbl{eq:sumegen}
\end{feqa}
with $E_b^\mathrm{max,min}$ of \ceqr{eq:ebeamlimits} and $E_{\gamma i}^\mathrm{max}=E_\gamma^\mathrm{tot,max}$ for each
enabled tail ($E_\gamma^\mathrm{tot,max}=0$ without radiation, \code{Egamma\_gen\_max} when set). The hadron range only
serves the diagnostic histograms.

\subsection{Generation volume}
\begin{feq}{reactions/EepReactionBase.cpp: EepReactionBase::genVolume}{simc.f:383-397}
V_\mathrm{gen} = \Delta\Omega_e\,\Delta\Omega_h\,\bigl(E'^\mathrm{gen}_\mathrm{max}-E'^\mathrm{gen}_\mathrm{min}\bigr)\quad[\mathrm{sr^2\,MeV}],
\lbl{eq:genvol}
\end{feq}
with $\Delta\Omega$ the products of the slope ranges (\ceqr{eq:genvol}). The Fermi momentum, the missing energy and the
$\Delta$ mass are drawn from normalised distributions and add nothing to $V_\mathrm{gen}$.

% =====================================================================================================================
\section{Completing the event}\lbl{sec:complete}
\subsection{Recoil mass of the \texorpdfstring{$\Delta$}{Delta} final states}
For \code{which\_pion} 2 and 3 every call of the event completion draws a new recoil mass from the table
\code{delta\_relativistic\_bw.inp} (a relativistic Breit--Wigner shape computed by P. Bosted, according to the code
comment; the first 1000 lines, masses $m_i$ from 1.028 to 1.627 GeV in 1 MeV steps, weights $b_i$; the file holds the
same table twice). With $B_i=\sum_{k\le i}b_k/\sum_k b_k$:
\begin{feq}{reactions/pion/PionProduction.cpp: DeltaMassTable::sample}{event.f:478-495, 620-631}
M_X = 1000\,m_{i^*}\ \mathrm{MeV},\qquad i^* = \text{the first } i \text{ with the smallest } |\xi-B_i|,
\lbl{eq:deltamass}
\end{feq}
i.e.\ the node nearest in cumulative probability, not an interpolated inverse. The distribution peaks at
$m=1.207$ GeV and 58\,\% of its weight lies below 1.232 GeV. A photon of tail 1 calls the completion again and draws
a new mass (\secr{sec:rad}).

\subsection{The pion energy}\lbl{sec:vertex}
With $\hat p$ the pion direction, $E_F$ and $\vec p_F=p_F\hat p_F$ of \eqr{eq:efer}, energy and momentum conservation
$\nu+E_F = E_h+E_X$, $\vec q+\vec p_F=\vec p_h+\vec p_X$, $E_X^2=M_X^2+p_X^2$ give, after squaring,
\begin{feqa}{reactions/pion/PionProduction.cpp: PionProduction::completeEv}{event.f:634-680}
a &= -|\vec q|\,(\hat q\cdot\hat p) - p_F\,(\hat p_F\cdot\hat p),\qquad
b = |\vec q|^2 + p_F^2 + 2|\vec q|\,p_F\,(\hat p_F\cdot\hat q),\qquad c = \nu + E_F,\nonumber\\
\tau &= c^2-b+M_h^2-M_X^2,\qquad
Q_A = 4(a^2-c^2),\quad Q_B = 4c\,\tau,\quad Q_C = -4a^2M_h^2-\tau^2,\nonumber\\
E_h &= \frac{-Q_B-\sqrt{Q_B^2-4Q_AQ_C}}{2Q_A},\qquad P_h=\sqrt{E_h^2-M_h^2},\qquad \delta_h = 100\,(P_h-P_h^0)/P_h^0,
\lbl{eq:quad}
\end{feqa}
where for H and coherent production $p_F=0$ and $c=\nu+M_A$ ($M_A=M_N$ there). The quadratic is that of
$\tau-2cE_h = 2aP_h$ squared. The try fails when the discriminant is negative, when $E_h<0$, when the recoil energy
$c-E_h\le M_X$ or when $E_h\le M_h$. Since $Q_A<0$, the chosen root is the larger energy, the forward solution in the
centre of mass; the code comment says that the second solution is hoped never to fall in the acceptance, and it prints
a warning in the first 5000 tries when that root lies inside the hadron energy edges. The sign condition
$\tau-2cE_h=2aP_h$ is not tested (\secr{sec:origin}). The vertex missing quantities follow from
\ceqr{eq:pm}; in addition $P_m=p_F$ (the generated value, not $|\vec P_m|$), $M_\mathrm{rec}=M_A-M_N+E_m$ and
\begin{feq}{reactions/pion/PionProduction.cpp: PionProduction::completeEv}{event.f:947-951}
T_\mathrm{rec} = 0\ \ (\text{H, coherent}),\qquad T_\mathrm{rec}=\sqrt{M_\mathrm{rec}^2+p_F^2}-M_\mathrm{rec}\ \ (\text{D},\ A\ge3).
\lbl{eq:trec}
\end{feq}

\subsection{Meson variables at the vertex}
\begin{feqa}{reactions/pion/PionProduction.cpp: PionProduction::pionMainKinematics}{event.f:723-771}
W^2 &= M_N^2+2M_N\nu-Q^2,\qquad \varepsilon = \Bigl[1+2\Bigl(1+\frac{\nu^2}{Q^2}\Bigr)\tan^2\frac{\theta_e}{2}\Bigr]^{-1},\qquad
\theta_{pq} = \arccos(\hat p\cdot\hat q),\nonumber\\
-t &= Q^2-M_h^2+2\nu E_h-2P_h|\vec q|\cos\theta_{pq},\qquad -t_\mathrm{min} = Q^2-M_h^2+2\nu E_h-2P_h|\vec q|,
\lbl{eq:tvertex}
\end{feqa}
and $\varphi_{pq}$ of \ceqr{eq:phipq}. The code keeps $-t$, the positive quantity $-(q-p_h)^2$, under the name
\code{t}. These values ignore the Fermi motion ($W$ is that of a nucleon at rest); the model computes its own
invariant mass with the Fermi motion (\eqr{eq:peepi}). For a deuteron or a heavier nucleus a relative momentum of the
recoil pair is stored (D: $k=\sqrt{\max(0,M^2-4M_X^2)}$ with $M$ the vertex missing mass; $A\ge3$:
$|\vec P_m+1.5\,\vec p_F|$, negative when $E_m<6$ MeV), which is not a tree column.

\subsection{\texorpdfstring{$\pi^0\to\gamma\gamma$}{pi0 -> 2 gamma}}
For \code{doing\_pizero} the completion decays the $\pi^0$ isotropically in its rest frame (two uniform numbers, first
the azimuth) and boosts the photons back along $\hat p$:
\begin{feqa}{physics/pion/PizeroDecay.cpp: pizeroDecay}{pizero\_decay.f:37-70; loren.f:18-23}
\beta &= \frac{P_h}{\sqrt{P_h^2+M_h^2}},\quad \gamma=(1-\beta^2)^{-1/2},\quad \varphi^\star=2\pi\xi_1,\quad \cos\theta^\star=2\xi_2-1,\quad
\vec k_{1}^\star = \tfrac{M_h}{2}\hat n(\theta^\star,\varphi^\star) = -\vec k_2^\star,\nonumber\\
E_i &= \gamma\bigl(\tfrac{M_h}{2}-\vec b\cdot\vec k_i^\star\bigr),\qquad
\vec k_i = \vec k_i^\star+\frac{\gamma^2}{1+\gamma}\,(\vec b\cdot\vec k_i^\star)\,\vec b-\gamma\,\tfrac{M_h}{2}\,\vec b,\qquad \vec b = -\beta\hat p,
\lbl{eq:pi0decay}
\end{feqa}
the photon four-momenta $(E_i,\vec k_i)$ in MeV in the laboratory frame being written to the tree. The code warns when
energy or momentum is not conserved to 0.001 MeV. The photons are traced to the calorimeter in \secr{sec:decay}.

\subsection{Reconstruction}\lbl{sec:recon}
From the reconstructed electron and pion (\ceqr{eq:recon}), with the hadron-arm reconstruction taken directly from the
spectrometer-entrance values for a $\pi^0$ (no $\pi^0$ reconstruction):
\begin{feqa}{reactions/EepReactionBase.cpp: EepReactionBase::completeRecon}{event.f:1325-1335, 1353-1354}
E_m^\mathrm{rec} &= \nu^\mathrm{rec}+M_N-E_h^\mathrm{rec},\qquad
m_N^2 = (E_m^\mathrm{rec})^2-(P_m^\mathrm{rec})^2,\qquad
m_A^2 = (\nu^\mathrm{rec}+M_A-E_h^\mathrm{rec})^2-(P_m^\mathrm{rec})^2,\nonumber\\
m &= \operatorname{sign}(m^2)\sqrt{|m^2|}\quad(\text{for } m_N, m_A),\qquad
-t^\mathrm{rec} = (Q^\mathrm{rec})^2-M_h^2+2\bigl(\nu^\mathrm{rec}E_h^\mathrm{rec}-P_h^\mathrm{rec}|\vec q^\mathrm{rec}|\cos\theta_{pq}^\mathrm{rec}\bigr),
\lbl{eq:recmm}
\end{feqa}
the tree columns \code{missmass} ($m_N$), \code{mmnuc} ($m_A$) and \code{t}. For a free nucleon $m_N$ peaks at $M_X$
(the neutron for $\pi^+$, the $\Delta$ distribution for \code{which\_pion} 2/3); the Fermi motion broadens it. The
reconstructed $W$ uses $M_p$ (\ceqr{eq:w}).

% =====================================================================================================================
\section{Jacobians}\lbl{sec:jac}
\subsection{Slopes}
Both slope pairs are generated, so the event jacobian is the product of the two solid-angle factors of \ceqr{eq:jacomega}:
\begin{feq}{reactions/pion/PionProduction.cpp: PionProduction::completeEv}{event.f:498, 1015-1025}
J = \bigl(1+x_e'^2+y_e'^2\bigr)^{-3/2}\bigl(1+x_h'^2+y_h'^2\bigr)^{-3/2}.
\lbl{eq:jac}
\end{feq}
No jacobian for the electron energy is needed: the pion energy is solved, and the cross section is differential in
$E'$, $\Omega_e$ and $\Omega_\pi$.

\subsection{Centre-of-mass frame and \texorpdfstring{$\mathrm{d}t\,\mathrm{d}\varphi^*\to\mathrm{d}\Omega_\pi$}{dt dphi -> dOmega}}\lbl{sec:cm}
The model gives $\mathrm{d}^2\sigma/\mathrm{d}t\,\mathrm{d}\varphi^*$ in the photon--nucleon centre of mass; the
transformation to the laboratory pion solid angle is computed by \code{transformToCm}, called by the model. Work in the
frame of \ceqr{eq:phipq} ($\hat z'\parallel\hat q$, $\hat y'\parallel\hat q\times\hat z$, $\hat x'=\hat y'\times\hat q$, in
the replay orientation of note 00). The pion has polar angle $\theta_{pq}$ and azimuth $\varphi_{pq}$, the Fermi momentum
the azimuth $\varphi_F$ (0 when its $x'$ component vanishes) and $\cos\theta_F'=\hat p_F\cdot\hat q$. The boost and the
Lorentz transformation of a four-vector $(E,\vec p)$ are
\begin{feqa}{reactions/EventChain.cpp: transformToCm}{jacobians.f:57-133; loren.f:18-23}
\vec\beta^* &= \frac{\bigl(p_Fp_{F,x'},\ p_Fp_{F,y'},\ |\vec q|+p_F\cos\theta_F'\bigr)}{\nu+E_F},\qquad
\gamma^* = (1-\beta^{*2})^{-1/2},\nonumber\\
E^* &= \gamma^*(E-\vec\beta^*\cdot\vec p),\qquad
\vec p^{\,*} = \vec p+\frac{\gamma^{*2}}{1+\gamma^*}(\vec\beta^*\cdot\vec p)\,\vec\beta^*-\gamma^*E\,\vec\beta^*,\nonumber\\
\cos\theta^* &= \frac{\vec p^{\,*}_h\cdot\vec q^{\,*}}{|\vec p^{\,*}_h|\,|\vec q^{\,*}|},
\lbl{eq:boost}
\end{feqa}
applied to the beam $(E_\mathrm{in}, E_\mathrm{in}\hat z)$, the photon $(\nu,|\vec q|\hat z')$, the pion
$(E_h, P_h\hat n(\theta_{pq},\varphi_{pq}))$ ($\hat n$ the unit vector of polar angle $\theta_{pq}$ and azimuth
$\varphi_{pq}$ in the frame of $\vec q$) and the nucleon $(E_F,\vec p_F)$ (a warning is printed when $|\beta^*|>1$). The azimuth $\varphi^*$ is that of $\vec p^{\,*}_h$ about
$\vec q^{\,*}$ in the axes $\hat y^*\parallel\vec q^{\,*}\times\vec k^{\,*}_\mathrm{beam}$, $\hat x^*\parallel\hat
y^*\times\vec q^{\,*}$, in $[0,2\pi)$. Without Fermi motion the boost is along $\vec q$ and $\varphi^*=\varphi_{pq}$.

The jacobian $|\partial(t,\varphi^*)/\partial(\cos\theta_{pq},\varphi_{pq})|$ is evaluated exactly for the non-collinear
boost. With $R=|\vec q|+p_F\cos\theta_F'-P_h\cos\theta_{pq}$, $\psi=\cos\varphi_F\cos\varphi_{pq}+\sin\varphi_F\sin\varphi_{pq}$
and $D=\nu+E_F-E_h$:
\begin{feqa}{reactions/EventChain.cpp: transformToCm}{jacobians.f:208-226}
\frac{\partial P_h}{\partial\cos\theta_{pq}} &=
\frac{P_h+\bigl(P_h^2\cos\theta_{pq}-\psi\,p_FP_h\sin\theta_F'\cos\theta_{pq}/\sin\theta_{pq}\bigr)/R}
{\bigl(DP_h/E_h+P_h\sin^2\theta_{pq}-\psi\,p_F\sin\theta_F'\sin\theta_{pq}\bigr)/R-\cos\theta_{pq}},\nonumber\\
\frac{\partial P_h}{\partial\varphi_{pq}} &=
\frac{p_FP_h\sin\theta_{pq}\sin\theta_F'\,\sin(\varphi_{pq}-\varphi_F)/R}
{\cos\theta_{pq}+\bigl(p_F\sin\theta_{pq}\sin\theta_F'\,\psi-P_h\sin^2\theta_{pq}-DP_h/E_h\bigr)/R},\nonumber\\
\frac{\partial t}{\partial\cos\theta_{pq}} &= 2\Bigl(|\vec q|P_h+\bigl(|\vec q|\cos\theta_{pq}-\nu\tfrac{P_h}{E_h}\bigr)\frac{\partial P_h}{\partial\cos\theta_{pq}}\Bigr),\qquad
\frac{\partial t}{\partial\varphi_{pq}} = 2\bigl(|\vec q|\cos\theta_{pq}-\nu\tfrac{P_h}{E_h}\bigr)\frac{\partial P_h}{\partial\varphi_{pq}},
\lbl{eq:dpdcos}
\end{feqa}
($t$ standing for the positive $-t$ of the code). For $\xi=\cos\theta_{pq}$ or $\varphi_{pq}$ the derivative of the boosted
pion momentum follows from \eqr{eq:boost} with $\partial E_h/\partial\xi=(P_h/E_h)\,\partial P_h/\partial\xi$, and in the
axes $(\hat x^*,\hat y^*)$
\begin{feqa}{reactions/EventChain.cpp: transformToCm}{jacobians.f:231-271}
\frac{\partial\vec p^{\,*}_h}{\partial\xi} &= \frac{\partial\vec p_h}{\partial\xi}
+\frac{\gamma^*-1}{\beta^{*2}}\Bigl(\vec\beta^*\cdot\frac{\partial\vec p_h}{\partial\xi}\Bigr)\vec\beta^*
-\gamma^*\vec\beta^*\,\frac{P_h}{E_h}\frac{\partial P_h}{\partial\xi},\qquad
\frac{\partial\varphi^*}{\partial\xi} = \frac{p^*_x\,\partial p^*_y/\partial\xi-p^*_y\,\partial p^*_x/\partial\xi}{p_x^{*2}+p_y^{*2}},\nonumber\\
J_\mathrm{cm} &= \Bigl|\frac{\partial t}{\partial\cos\theta_{pq}}\frac{\partial\varphi^*}{\partial\varphi_{pq}}
-\frac{\partial t}{\partial\varphi_{pq}}\frac{\partial\varphi^*}{\partial\cos\theta_{pq}}\Bigr|\quad[\mathrm{MeV^2/sr}],
\lbl{eq:jcm}
\end{feqa}
the code writing the components of $\partial\vec p_h/\partial\xi$ out term by term. $J_\mathrm{cm}$ enters the cross
section (\eqr{eq:peepi}), not $J$. Without Fermi motion it reduces to $\partial t/\partial\cos\theta_{pq}$, the
$2|\vec q|P_h\cdots$ factor of the collinear case. The code also keeps the older collinear approximation
\begin{feq}{reactions/EventChain.cpp: transformToCm}{jacobians.f:275-278}
J_\mathrm{old} = \frac{2\bigl(E_F-2p_F\hat p_{F,z}\tfrac{E_h}{P_h}\cos\theta_{pq}\bigr)\bigl(|\vec q|+p_F\hat p_{F,z}\bigr)P_h}
{E_F+\nu-\bigl(|\vec q|+p_F\hat p_{F,z}\bigr)\tfrac{E_h}{P_h}\cos\theta_{pq}}-2P_hp_F,
\lbl{eq:jold}
\end{feq}
which is stored in the main record and not used ($\hat p_{F,z}$ is the component along the beam, not along $\vec q$).

% =====================================================================================================================
\section{Cross-section model}\lbl{sec:xsec}
\subsection{What the reaction requires from a model}\lbl{sec:require}
\code{PionProduction::completeMain} sets $S=1$ and calls \code{sigccMain(vertex, main)} of the run's
\code{IPhysicsModel}. The model reads the vertex record ($E_\mathrm{in}$, $E'$, $\nu$, $Q^2$, $\vec q$, the pion), the main
record ($\varepsilon$, $-t$, $\varphi_{pq}$, the Fermi momentum $p_F$, $\hat p_F$, $E_F$) and the set-up values $M_N$,
\code{which\_pion} and \code{doing\_pizero}. It must return the cross section per nucleon
$\mathrm{d}^5\sigma/\mathrm{d}E'\mathrm{d}\Omega_e\mathrm{d}\Omega_\pi$ in $\mu$b/(MeV sr$^2$), with the pion solid angle
in the laboratory, and may write the centre-of-mass angles, $W^*$, the jacobians and the tree value \code{sigcm} into
the main record; \code{sigccRecon} returns 1. The reaction then applies the factors of \eqr{eq:cgfactors}, the
Coulomb factor and $n_\mathrm{tgt}$.

\subsection{The model \texttt{pion}: from the centre of mass to the laboratory}
\begin{feqa}{physics/pion/PionCrossSections.cpp: peepi}{physics\_pion.f:74-92, 126-191}
s &= (\nu+E_F)^2-\bigl(|\vec q|+p_F\cos\theta_F'\bigr)^2-\bigl(p_F\sin\theta_F'\bigr)^2,\qquad W^*=\sqrt s,\nonumber\\
k_\mathrm{eq} &= \frac{W^{*2}-M_N^2}{2M_N},\qquad
\Gamma = \frac{\alpha}{2\pi^2}\,\frac{E'}{E_\mathrm{in}}\,\frac{k_\mathrm{eq}}{Q^2}\,\frac{1}{1-\varepsilon},\qquad
f = \frac{1}{1-p_F\hat p_{F,z}/E_F}\,\frac{M_N}{E_F},\nonumber\\
\sigma &= \frac{\mathrm{d}^2\sigma}{\mathrm{d}t\,\mathrm{d}\varphi^*}\cdot J_\mathrm{cm}\cdot\Gamma\cdot f
\quad[\mu\mathrm{b/(MeV\,sr^2)}],
\lbl{eq:peepi}
\end{feqa}
($\cos\theta_F'$ set to 1 when it exceeds 1 by less than $10^{-8}$). $\Gamma$ is the Hand virtual-photon flux with the
equivalent photon energy $k_\mathrm{eq}$ of a nucleon at rest; the code comments explain the choice of $M_N$ instead of
the Fermi-moving nucleon energy (a 2000 version used $E_F-p_F\cos\theta_F'$, the 2001 version returned to $M_N$
``more consistent with the phase space''), and $f$ as the non-collinear flux correction: the relative velocity
$1-p_{F,z}/E_F$ of electron and nucleon and $E_F$ instead of $M_N$ in the flux $|v_1-v_2|\,2E_1\,2E_2$.
$\mathrm{d}^2\sigma/\mathrm{d}t\,\mathrm{d}\varphi^*$ (in $\mu$b/(MeV$^2$ rad)) is the blend of \eqr{eq:blend}.

\subsection{The 2021 fit of P. Bosted}
The default centre-of-mass cross section is a fit to exclusive $\pi^+$ and $\pi^-$ data (code comment: ``April 2021 fit
to exclusive pi+ and pi- data from fpi1, fpi2, CT, pt-SIDIS, CSV-SIDIS, and KLT''), evaluated with $s$, $Q^2$ and $-t$ in
GeV$^2$ and the centre-of-mass angles of \secr{sec:cm}. With the parameter set $p_1\ldots p_{17}$ of \tabr{tab:pfit}:
\begin{feqa}{physics/pion/PionCrossSections.cpp: exclfit}{physics\_pion.f:807-851}
F_\pi &= \frac{\phi}{1+p_1Q^2+p_2Q^4},\qquad
\sigma_L = \Bigl(p_3+\frac{p_{15}}{Q^2}\Bigr)\frac{|t|}{(|t|+0.02)^2}\,Q^2F_\pi^2\,e^{p_4|t|}\Big/\bigl(s^{p_{11}}+s^{p_{17}/2}\bigr),\nonumber\\
\sigma_T &= \frac{p_5}{Q^2}\,e^{p_6Q^4}\,e^{p_{14}|t|}\Big/\bigl(s^{p_{12}}+s^{p_{16}/2}\bigr),\qquad
\sigma_{LT} = \frac{p_7}{1+p_{10}Q^2}\,e^{p_8|t|}\,\frac{\sin\theta^*}{s^{p_{13}}},\nonumber\\
\sigma_{TT} &= \frac{p_9}{1+Q^2}\,e^{-7|t|}\sin^2\theta^*,\qquad
\sigma_\mathrm{u} = \sigma_T+\varepsilon\sigma_L+\varepsilon\cos2\varphi^*\,\sigma_{TT}+\sqrt{2\varepsilon(1+\varepsilon)}\cos\varphi^*\,\sigma_{LT},\nonumber\\
\frac{\mathrm{d}^2\sigma}{\mathrm{d}t\,\mathrm{d}\varphi^*}\Big|_{2021} &= \frac{8.539\,\sigma_\mathrm{u}}{(s-0.938^2)^2}\,\frac{1}{2\cdot3.1415928\cdot10^6},
\lbl{eq:exclfit}
\end{feqa}
(the last factor converts GeV$^{-2}$ to MeV$^{-2}$ and $\mathrm{d}t\to\mathrm{d}t\,\mathrm{d}\varphi^*$). The charged
pion uses $\phi=1$ and the set $p^+$ for \code{which\_pion} 0, 2, 10, $p^-$ for 1, 3, 11; the $\pi^0$ uses $\phi=0$
(no pion pole, $\sigma_L=0$) and the average of the results of the two sets:
\begin{feq}{physics/pion/PionCrossSections.cpp: sigParam2021}{physics\_pion.f:788-803}
\frac{\mathrm{d}^2\sigma}{\mathrm{d}t\,\mathrm{d}\varphi^*}\Big|_{2021,\pi^0} = \tfrac12\Bigl(
\frac{\mathrm{d}^2\sigma}{\mathrm{d}t\,\mathrm{d}\varphi^*}\Big|_{p^+,\phi=0}+\frac{\mathrm{d}^2\sigma}{\mathrm{d}t\,\mathrm{d}\varphi^*}\Big|_{p^-,\phi=0}\Bigr).
\lbl{eq:pizerofit}
\end{feq}
Remarks from the code: the normalisation 8.539 and the denominator $(s-M^2)^2$ come from earlier fits scaled to a fixed
$W$; the comment says $W=1.96$ GeV, while $(W^2-0.938^2)^2=8.539$ corresponds to $W=1.950$ GeV (with 1.96 it would be
8.772). The proton mass is fixed at 0.938 GeV inside the fit and $\pi$ is written 3.1415928. The fit is used as it is
for coherent production ($s$ of the photon--nucleus system); it is not clamped and can be negative. The parameters are
transcribed from the code; there is no publication to compare them with (the code cites the output of a private
program, \code{ptb.f}).

\begin{simctable}{footnotesize}{L{0.12\linewidth}L{0.2\linewidth}L{0.2\linewidth}L{0.48\linewidth}}
\caption{Parameters of the 2021 fit \fsrc{physics\_pion.f:751-786} (one parameter per row).}\lbl{tab:pfit}\\
\toprule\rowcolor{simchead}\simcth{parameter} & \simcth{$p^+$ ($\pi^+$)} & \simcth{$p^-$ ($\pi^-$)} & \simcth{role in \eqr{eq:exclfit}}\\\midrule\endfirsthead
\toprule\rowcolor{simchead}\simcth{parameter} & \simcth{$p^+$ ($\pi^+$)} & \simcth{$p^-$ ($\pi^-$)} & \simcth{role in \eqr{eq:exclfit}}\\\midrule\endhead
\bottomrule\endlastfoot
$p_1$ & 1.60077 & 1.75169 & $Q^2$ term of the form factor\\
$p_2$ & $-0.01523$ & 0.11144 & $Q^4$ term of the form factor\\
$p_3$ & 37.08142 & 47.35877 & $\sigma_L$ normalisation\\
$p_4$ & $-4.11060$ & $-4.69434$ & $\sigma_L$ $t$ slope\\
$p_5$ & 23.26192 & 1.60552 & $\sigma_T$ normalisation\\
$p_6$ & 0.00983 & 0.00800 & $\sigma_T$ $Q^4$ exponent\\
$p_7$ & 0.87073 & 0.44194 & $\sigma_{LT}$ normalisation\\
$p_8$ & $-5.77115$ & $-2.29188$ & $\sigma_{LT}$ $t$ slope\\
$p_9$ & $-271.08678$ & $-41.67194$ & $\sigma_{TT}$ normalisation\\
$p_{10}$ & 0.13766 & 0.69475 & $\sigma_{LT}$ $Q^2$ term\\
$p_{11}$ & $-0.00855$ & 0.02527 & $\sigma_L$ $s$ power\\
$p_{12}$ & 0.27885 & $-0.50178$ & $\sigma_T$ $s$ power\\
$p_{13}$ & $-1.13212$ & $-1.22825$ & $\sigma_{LT}$ $s$ power\\
$p_{14}$ & $-1.50415$ & $-1.16878$ & $\sigma_T$ $t$ slope\\
$p_{15}$ & $-6.34766$ & 5.75825 & $\sigma_L$ $1/Q^2$ term\\
$p_{16}$ & 0.55769 & $-1.00355$ & $\sigma_T$ $\sqrt s$ power\\
$p_{17}$ & $-0.01709$ & 0.05055 & $\sigma_L$ $\sqrt s$ power\\
\end{simctable}

\subsection{MAID2007 below \texorpdfstring{$W=2$}{W = 2} GeV}
For $W^*<2$ GeV the code also evaluates the unpolarised MAID2007 cross section from tables produced by P. Bosted with
code of L. Tiator (code comment), four files of 25 $Q^2$ $\times$ 46 $W$ $\times$ 23 angle rows of 18 response
functions: \code{maidpi0p.dat} ($\pi^0p$, $i_\pi=1$), \code{maidpi0n.dat} ($\pi^0n$, 2), \code{maidpipn.dat}
($\pi^+n$, 3), \code{maidpimp.dat} ($\pi^-p$, 4). $\pi^0$ uses $i_\pi=2$ for \code{which\_pion} 1 and 3, 1 otherwise;
the charged pion uses 4 for \code{which\_pion} 1, 3, 11 and 3 otherwise. With $Q^2$, $W$, $E_0=E_\mathrm{in}$ in GeV,
$m=0.9383$:
\begin{feqa}{physics/pion/PionCrossSections.cpp: sigmaid}{physics\_pion.f:636-705}
\nu_0 &= \frac{W^2-m^2+Q^2}{2m},\qquad \sin^2\tfrac{\theta_0}{2}=\frac{Q^2}{4E_0(E_0-\nu_0)},\nonumber\\
\varepsilon_0 &= \Bigl[1+2\Bigl(1+\frac{\nu_0^2}{Q^2}\Bigr)\frac{\sin^2(\theta_0/2)}{1-\sin^2(\theta_0/2)}\Bigr]^{-1},\nonumber\\
i_Q &= \min\bigl(25,\max(1,\lfloor(Q^2+0.1)/0.2\rfloor)\bigr),\qquad
i_W = \min\bigl(46,\max(1,\lfloor(W-1.09)/0.02\rfloor)\bigr),\nonumber\\
w_f &= 1\ (W\le1.232),\qquad w_f=(W-1.132)/0.1\ (W>1.232),\nonumber\\
\sigma_T &= \frac{R_1}{\max(0.2,Q^2)\,w_f},\qquad \sigma_L=R_2\sigma_T,\qquad \sigma_{LT}=R_3\sigma_T,\qquad \sigma_{TT}=R_4\sigma_T,\nonumber\\
\frac{\mathrm{d}\sigma}{\mathrm{d}\Omega^*}\Big|_\mathrm{MAID} &= \sigma_T+\varepsilon_0\sigma_L\nonumber\\
&\quad+\sqrt{2\varepsilon_0(1+\varepsilon_0)}\cos\varphi^*\,\sigma_{LT}+\varepsilon_0\cos2\varphi^*\,\sigma_{TT}\quad[\mu\mathrm{b/sr}],
\lbl{eq:maid}
\end{feqa}
with $R_j=R_j(i_\pi,i_Q,i_W,i_\theta)$ the table entries and $i_\theta$ the bin of $\cos\theta^*$ in the edges
$-0.20$, 0.20, 0.44, 0.63, 0.78, 0.90, 1.0 (the last bin containing it; $\cos\theta^*<-0.2$ uses bin 1). The result is 0
for $W<1.08$ GeV, $\nu_0>E_0$, or $\sin^2(\theta_0/2)\notin(0,1]$. The table is used at the bin found by truncation,
without interpolation, and only the first six angle rows are reachable. $\varepsilon_0$ is that of a nucleon at rest at
$(W,Q^2,E_0)$, not the event's $\varepsilon$. The target-polarisation responses $R_5\ldots R_{18}$ are combined into
polarised cross sections that are not used. MAID2007 is the model of the reference cited below; the tables were not
compared with it for this note.

The model value, in $\mu$b/(MeV$^2$ rad), is the blend
\begin{feq}{physics/pion/PionCrossSections.cpp: peepi}{physics\_pion.f:133-154}
\frac{\mathrm{d}^2\sigma}{\mathrm{d}t\,\mathrm{d}\varphi^*} =
\begin{cases} f_W\,\sigma_{2021}+(1-f_W)\,\dfrac{\mathrm{d}\sigma/\mathrm{d}\Omega^*|_\mathrm{MAID}}{2\,p^*_h\,q^*}, & W^*<2\ \mathrm{GeV},\\
\sigma_{2021}, & W^*\ge2\ \mathrm{GeV},\end{cases}\qquad f_W = \min\Bigl(1,\max\bigl(0,\tfrac{W^*-1.5}{0.4}\bigr)\Bigr),
\lbl{eq:blend}
\end{feq}
with $\sigma_{2021}$ of \eqsr{eq:exclfit}{eq:pizerofit}, $p^*_h$, $q^*$ the centre-of-mass momenta in MeV
($\mathrm{d}t=2p^*q^*\,\mathrm{d}\cos\theta^*$), $W^*$ in GeV in $f_W$. Below 1.5 GeV the cross section is MAID, above
1.9 GeV the fit. The tree column \code{sigcm} is this value.

\subsection{Fits present in the code and not used}
\code{sig\_blok} (a fit to the data of Brauel et al.\ \citer{brauel}, used by the Delta reaction, note 05), \code{sig\_param04}
and \code{sig\_param\_3000} (earlier fits to the $F_\pi$ data) and \code{sig\_multipole} (a low-$W$ multipole model with
CGLN \citer{cgln}, Walker \citer{walker} and Bartl \citer{bartl} amplitudes, \code{MULTIPOL.file}) are ported but their
calls are commented out in the Fortran \fsrc{physics\_pion.f:94-123, 156-163}; SIMC++ keeps them for the L1 tests.

\subsection{\texorpdfstring{$\Delta$}{Delta} final states and the number of target nucleons}
The cross section of $\gamma^*N\to\pi\Delta$ is the nucleon-recoil cross section scaled by empirical factors (code
comment: ``Clebsch-Gordon coefficients ... updated from Peter Bosted's empirical checks'', 2021 and 2023). SIMC++
applies, in both modes,
\begin{feqa}{reactions/pion/PionProduction.cpp: PionProduction::completeMain}{event.f:1461-1503 (294f662: event.f:1458-1500)}
\text{\code{which\_pion} 2:}&\quad \sigma\leftarrow 0.55\,\sigma\ (\pi^0,\ \text{H, D}),\quad
\sigma\leftarrow0.4\,\sigma\ (\pi^+,\ \text{H}),\quad \sigma\leftarrow(0.4+0.8)\,\sigma\ (\pi^+,\ \text{D}),\nonumber\\
\text{\code{which\_pion} 3:}&\quad \sigma\leftarrow0.99\,\sigma\ (\pi^0,\ \text{D}),\quad
\sigma\leftarrow0.55\,\sigma\ (\pi^-,\ \text{H}),\quad \sigma\leftarrow(0.55+0.99)\,\sigma\ (\pi^-,\ \text{D}),\nonumber\\
n_\mathrm{tgt} &= N\ (\text{\code{which\_pion} 1, 11}),\qquad n_\mathrm{tgt}=Z\ (\text{all other values}),
\lbl{eq:cgfactors}
\end{feqa}
and no factor for $A\ge3$ (the branches exist for H and D only), so for $A\ge3$ the $\Delta$ final states are weighted
with $Z$ times the nucleon-recoil cross section. Upstream SIMC 59ccdf6 changed this block: the D factors became
$Z\cdot0.4+N\cdot0.8$ and $Z\cdot0.55+N\cdot0.99$ (and $Z$ times the H factors), and $n_\mathrm{tgt}$ is set only for
\code{which\_pion} 0, 10 ($Z$) and 1, 11 ($N$), so that it stays 1 for 2 and 3. For H and D the weight is unchanged
($Z=N=1$); for $A\ge3$ and \code{which\_pion} 2/3 the upstream weight is the SIMC++ weight divided by $Z$. SIMC++ keeps
the earlier weight in both modes until the author of the change has answered (\secr{sec:origin}). For hydrogen
$\gamma^*n\to\pi^0\Delta^0$ is set to 0 (and refused at set-up).

\subsection{Weight and units}
\begin{feq}{reactions/pion/PionProduction.cpp: PionProduction::completeMain}{event.f:1536-1538, 1561-1562}
w = J\cdot w_\mathrm{gen}\cdot\sigma\cdot n_\mathrm{tgt}\quad[\mu\mathrm{b/(MeV\,sr^2)}],\qquad
\sigma\leftarrow\sigma\,(1+V_C/E_b)^2\ \text{with the Coulomb correction},
\lbl{eq:weight}
\end{feq}
with $\sigma$ of \eqr{eq:peepi} times the factors of \eqr{eq:cgfactors}, $J$ of \eqr{eq:jac} and $w_\mathrm{gen}$ of
\eqr{eq:genee} times the radiative weight. With $V_\mathrm{gen}$ in sr$^2$ MeV the products $w\,\mathcal L\,V_\mathrm{gen}
/N_\mathrm{tried}$ are counts (\ceqr{eq:normfac}). The \texttt{.hist} line \code{CENTRAL.sigcc} is the model at the
central kinematics; in the Fortran (and the compat mode) the central event has $\varepsilon=-t=\varphi_{pq}=0$, while the
default mode sets the central $W$, $\varepsilon$, $\theta_{pq}$, $-t$ and $\varphi_{pq}$ first
(\code{PionProduction::prepareCentral}).

% =====================================================================================================================
\section{Radiation, energy loss and decays}\lbl{sec:rad}
\subsection{Radiative tails}
All three legs radiate; the pion tail uses $\lambda_3$ of \ceqr{eq:lambda} with the pion energy and momentum. The photon
limits of this reaction are
\begin{feqa}{radiation/Radc.cpp: generateRad}{radc.f:297-308, 372-389, 428-446}
E_{\gamma1} &\in \bigl[-d_E,\ \min(\Sigma_\mathrm{max}-E',\ E_{\gamma1}^\mathrm{max})+d_E\bigr],\nonumber\\
E_{\gamma2} &\in \bigl[E'-E'_\mathrm{max}-d_E,\ \min(E'-E'_\mathrm{min},\ E_\gamma^\mathrm{tot,max}-E_{\gamma1})+d_E\bigr],\nonumber\\
E_{\gamma3} &\in \bigl[E_h-E_{h,\mathrm{max}}-d_E,\ \min(E_h-E_{h,\mathrm{min}},\ E_\gamma^\mathrm{tot,max}-E_{\gamma1}-E_{\gamma2})+d_E\bigr],
\lbl{eq:tail1}
\end{feqa}
(\code{Egamma\_gen\_max} replaces the upper limits when set; the code comment calls the exact beam-energy limits of
tail 1 ``too complicated'' for pion production and uses the remaining generated energy). A tail-1 photon lowers
$E_\mathrm{in}$ and the event completion runs again with it: the pion energy is solved again and, for \code{which\_pion}
2/3, a new $\Delta$ mass is drawn and, for a $\pi^0$, a new photon pair. The soft-photon factors assume elastic e--p
kinematics (\ceqr{eq:bremkin}), and the internal radiator constant of the hadron uses $M_p$ for every hadron (note 00).

\subsection{Energy loss and multiple scattering}
The pion loses energy in the target and the windows with its own mass in the Bethe--Bloch formula and scatters with
$\beta=P_h/E_h$ (\csecr{sec:eloss}); the reconstruction adds back the most probable loss with \code{correct\_Eloss}.
With \code{using\_HMScoll} or \code{using\_SHMScoll} the spectrometer collimator is stepped through, with energy loss,
multiple scattering and a pion absorption probability; this belongs to the spectrometer models and is not described here.

\subsection{Decay in flight}\lbl{sec:decay}
With \code{doing\_decay} = 1 a charged pion can decay inside the HMS, SHMS, SOS and HRS models, along the whole path
from the target plane (the first drift of each model starts there) to the detector hut; the mean decay length uses the input \code{ctau} (cm, 780.45 cm for the pion
is the usual value; the code does not set it). In each field-free drift of length $L$ (along the central ray) the decay
point is drawn and tested against the path of the ray:
\begin{feqa}{optics/CosyEngine.cpp: project}{shared/project.f:45-69}
\beta &= \frac{p}{\sqrt{p^2+m^2}},\qquad \lambda = c\tau\,\beta\gamma,\qquad z_d=-\lambda\ln(1-\xi),\qquad
\ell = L\sqrt{1+x'^2+y'^2},\nonumber\\
z_d>\ell&:\ \ \text{no decay; } d\leftarrow d+\ell,\nonumber\\
z_d\le\ell&:\ \ \text{decay at } (x,y)\leftarrow(x,y)+(x',y')\,\frac{z_d}{\sqrt{1+x'^2+y'^2}},\ \ d\leftarrow d+z_d,
\lbl{eq:decaypoint}
\end{feqa}
$d$ being the tree column \code{decdist} (reset to 0 for every try). The decay is isotropic in the pion rest frame,
$\pi\to\mu\nu$ with the muon momentum $p^\star=29.783$ MeV/$c$ (the masses of note 00 give 29.792 MeV/$c$), and the muon
is boosted along the track with \eqr{eq:boost} ($\vec\beta=-\beta\,(x',y',1)/\sqrt{1+x'^2+y'^2}$ in the loren
convention):
\begin{feqa}{optics/CosyEngine.cpp: decayProducts}{shared/project.f:72-116; shared/transp.f:146-184}
\varphi^\star&=2\pi\xi_1,\quad \cos\theta^\star=2\xi_2-1,\quad \vec p^{\,\star}=p^\star\hat n(\theta^\star,\varphi^\star),\quad
E^\star=\sqrt{M_\mu^2+p^{\star2}},\nonumber\\
x' &\leftarrow p_x/p_z,\quad y'\leftarrow p_y/p_z,\quad \delta\leftarrow100\,(p/p_\mathrm{spec}-1),\quad m\leftarrow M_\mu,\nonumber\\
x &\leftarrow x+x'\,\Bigl(L-\frac{z_d}{\sqrt{1+x'^2+y'^2}}\Bigr)\ \ (\text{the rest of the drift}),
\lbl{eq:decayproducts}
\end{feqa}
where the Fortran computes the remaining drift with the muon slopes, the default mode with the pion slopes (so that
the track ends at the end of the drift). Inside a magnetic element of length $L$, whose transport is a COSY map that
cannot be split, the decay is applied at the entrance when $z_d\le L/2$ (and the map transports the muon) or at the
exit when $L/2<z_d\le L+\Delta L$ (with the pion's $\beta$, $\gamma$; $\Delta L$ the path-length correction of the map);
$d$ grows by $z_d$ or by $L+\Delta L$ \fsrc{shared/transp.f:132-185, 231-276}. After a decay the muon is transported,
scattered and reconstructed as if it were the pion, and \code{Mhadron} is $M_\mu$. Kaon decays use the same code
(note on kaon production). A $\pi^0$ never decays in flight in SIMC++; the Fortran stops at the first decay test of a
$\pi^0$ or of one of its photons (its mass is not within 2 MeV of $M_\pi$).

\subsection{\texorpdfstring{$\pi^0$}{pi0} photons in a calorimeter}
With \code{doing\_pizero} and a calorimeter as hadron arm (codes 7, 8) each photon of \eqr{eq:pi0decay} is rotated
into the frame of the calorimeter centred at the hadron-arm angle $\theta_h^0$ and traced to its face:
\begin{feqa}{driver/PizeroCalorimeter.cpp: rotateToCalorimeter}{simc.f:1506-1526}
\text{arm 8 (beam left):}&\quad k'_x=k_x,\quad k'_y=k_y\cos\theta_h^0-k_z\sin\theta_h^0,\quad k'_z=k_y\sin\theta_h^0+k_z\cos\theta_h^0,\nonumber\\
\text{arm 7 (beam right):}&\quad k'_x=k_x,\quad k'_y=k_y\cos\theta_h^0+k_z\sin\theta_h^0,\quad k'_z=-k_y\sin\theta_h^0+k_z\cos\theta_h^0,\nonumber\\
x'_\gamma &= k'_x/k'_z,\qquad y'_\gamma=k'_y/k'_z,
\lbl{eq:calorot}
\end{feqa}
from the vertex position of the hadron arm. The try succeeds if both photons hit (\code{pizero\_ngamma} = 2) or at least
one (\code{pizero\_ngamma} = 1). The hit positions go to the columns \code{xcal\_gamma1/2}, \code{ycal\_gamma1/2}
($-10^{10}$ for a missed photon) and the hadron-arm focal-plane columns hold the last photon that hit. The photons are
massless and stable in the default mode; each photon's path starts at 0. With a magnetic spectrometer as hadron arm
the $\pi^0$ itself is transported as if it were charged, and its photons are only written to the tree.

% =====================================================================================================================
\section{Inputs and outputs}\lbl{sec:io}
\subsection{Inputs}
\begin{simctable}{footnotesize}{L{0.24\linewidth}L{0.12\linewidth}L{0.08\linewidth}L{0.56\linewidth}}
\caption{Inputs specific to pion production (one input per row; the common inputs are in note 01, \S7, and the user
manual, chapter 4).}\lbl{tab:inputs}\\
\toprule\rowcolor{simchead}\simcth{input} & \simcth{unit} & \simcth{default} & \simcth{meaning here}\\\midrule\endfirsthead
\toprule\rowcolor{simchead}\simcth{input} & \simcth{unit} & \simcth{default} & \simcth{meaning here}\\\midrule\endhead
\bottomrule\endlastfoot
\code{doing\_pion} & flag & 0 & 1 selects this reaction\\
\code{which\_pion} & code & 0 & final state (\tabr{tab:finalstates})\\
\code{doing\_pizero} & flag & 0 & $\pi^0$ instead of a charged pion\\
\code{pizero\_ngamma} & 1 or 2 & 0 & photons required in the calorimeter (must be set for a $\pi^0$)\\
\code{drift\_to\_cal} & cm & 0 & distance from the target to the calorimeter face (arms 7, 8)\\
\code{doing\_decay} & flag & 0 & decay in flight of the charged pion\\
\code{ctau} & cm & 0 & $c\tau$ of the hadron for the decay\\
\code{using\_HMScoll} & flag & 0 & HMS collimator stepping and pion absorption\\
\code{using\_SHMScoll} & flag & 0 & SHMS collimator stepping and pion absorption\\
\code{targ\%A} & -- & 0 & H, D or $A\ge3$ case; Fermi-motion table\\
\code{targ\%Z} & -- & 0 & $n_\mathrm{tgt}$, Coulomb correction\\
\code{targ\%mass\_amu} & u & 0 & $M_A$ (forced for H and D)\\
\code{targ\%mrec\_amu} & u & 0 & $M_\mathrm{rec}$ for $A\ge3$; final nucleus for coherent production (replaced when off by 100 MeV)\\
\code{hadron\_arm} & code & 0 & pion spectrometer (7/8: calorimeter for $\pi^0$)\\
\code{spec\%p\%P} & MeV/$c$ & 0 & $P_h^0$\\
\code{spec\%p\%theta} & deg & 0 & $\theta_h^0$; calorimeter angle\\
\code{SPedge\%p\%yptar\%min} & mrad & 0 & generation range of $y'_h$, lower limit\\
\code{SPedge\%p\%yptar\%max} & mrad & 0 & generation range of $y'_h$, upper limit\\
\code{SPedge\%p\%xptar\%min} & mrad & 0 & generation range of $x'_h$, lower limit\\
\code{SPedge\%p\%xptar\%max} & mrad & 0 & generation range of $x'_h$, upper limit\\
\code{SPedge\%p\%delta\%min} & \% & 0 & hadron energy edges (limits, tail 3)\\
\code{SPedge\%p\%delta\%max} & \% & 0 & hadron energy edges (limits, tail 3)\\
\code{using\_tgt\_field} & flag & 0 & polarised-target field: tracking and the angles to the polarisation\\
\end{simctable}
The inputs \code{deForest\_flag}, \code{use\_benhar\_sf}, \code{transparency} and the $E_m$ cut have no effect on
this reaction.

\subsection{Outputs}
% docs/physics/latex/frag_pion_columns.tex - the meson tree layout of the pion and Delta reactions, \input by the notes
% 04 (pion_production) and 05 (delta_production). C++: output/NtupleLayout.cpp (UpstreamMesonColumns,
% FillUpstreamNtuple; legacy layout PionNtupleColumns, FillPionNtuple). Fortran: NtupleInit.f:34-224 and
% results_write.f:120-286 (upstream 59ccdf6). One column per row. No displayed equations here.
The sink \code{root\_tree} writes one entry of the tree \code{h10} per success. With the layout of upstream SIMC 59ccdf6
(the default) the tree has the 64 columns of \tabr{tab:columns}; \code{-{}-tree-layout legacy} gives the layout of SIMC
294f662. Columns 1--24 are those of the (e,e\pr p) layout: \code{hs} is the arm on the beam right, \code{ss} the arm on
the beam left, whatever particle each detects; the reconstructed columns come from the \code{recon} record, the
``true'' ones (suffix \code{i}) from the \code{orig} record for the arm columns and from the \code{vertex} record
for the physics columns. ``Hadron'' is the detected hadron: the pion in note 04, the proton in note 05.

\begin{simctable}{footnotesize}{L{0.06\linewidth}L{0.15\linewidth}L{0.13\linewidth}L{0.66\linewidth}}
\caption{Tree columns of the pion and Delta reactions (default layout; one column per row).}\lbl{tab:columns}\\
\toprule\rowcolor{simchead}\simcth{\#} & \simcth{name} & \simcth{unit} & \simcth{content}\\\midrule\endfirsthead
\toprule\rowcolor{simchead}\simcth{\#} & \simcth{name} & \simcth{unit} & \simcth{content}\\\midrule\endhead
\bottomrule\endlastfoot
1 & \code{hsdelta} & \% & reconstructed $\delta$, \code{hs} arm\\
2 & \code{hsyptar} & rad & reconstructed $y'$, \code{hs} arm\\
3 & \code{hsxptar} & rad & reconstructed $x'$, \code{hs} arm\\
4 & \code{hsytar} & cm & reconstructed horizontal target position, \code{hs} arm\\
5 & \code{hsxfp} & cm & focal-plane $x$, \code{hs} arm\\
6 & \code{hsxpfp} & rad & focal-plane $x'$, \code{hs} arm\\
7 & \code{hsyfp} & cm & focal-plane $y$, \code{hs} arm\\
8 & \code{hsypfp} & rad & focal-plane $y'$, \code{hs} arm\\
9 & \code{hsdeltai} & \% & true $\delta$ after radiation, \code{hs} arm\\
10 & \code{hsyptari} & rad & true $y'$, \code{hs} arm\\
11 & \code{hsxptari} & rad & true $x'$, \code{hs} arm\\
12 & \code{hsytari} & cm & $z_t\sin\theta^0$ of the \code{hs} arm\\
13 & \code{ssdelta} & \% & reconstructed $\delta$, \code{ss} arm\\
14 & \code{ssyptar} & rad & reconstructed $y'$, \code{ss} arm\\
15 & \code{ssxptar} & rad & reconstructed $x'$, \code{ss} arm\\
16 & \code{ssytar} & cm & reconstructed horizontal target position, \code{ss} arm\\
17 & \code{ssxfp} & cm & focal-plane $x$, \code{ss} arm\\
18 & \code{ssxpfp} & rad & focal-plane $x'$, \code{ss} arm\\
19 & \code{ssyfp} & cm & focal-plane $y$, \code{ss} arm\\
20 & \code{ssypfp} & rad & focal-plane $y'$, \code{ss} arm\\
21 & \code{ssdeltai} & \% & true $\delta$ after radiation, \code{ss} arm\\
22 & \code{ssyptari} & rad & true $y'$, \code{ss} arm\\
23 & \code{ssxptari} & rad & true $x'$, \code{ss} arm\\
24 & \code{ssytari} & cm & $-z_t\sin\theta^0$ of the \code{ss} arm\\
25 & \code{q} & GeV/$c$ & $|\vec q^\mathrm{rec}|$\\
26 & \code{nu} & GeV & $\nu^\mathrm{rec}$\\
27 & \code{Q2} & (GeV/$c$)$^2$ & $(Q^\mathrm{rec})^2$\\
28 & \code{W} & GeV & $W^\mathrm{rec}$, computed with the proton mass\\
29 & \code{xbj} & -- & $x^\mathrm{rec}$\\
30 & \code{epsilon} & -- & $\varepsilon^\mathrm{rec}$\\
31 & \code{Em} & GeV & $E_m^\mathrm{rec}=\nu^\mathrm{rec}+M_N-E_h^\mathrm{rec}$\\
32 & \code{Pm} & GeV/$c$ & $P_m^\mathrm{rec}=|\vec P_m^\mathrm{rec}|$\\
33 & \code{qi} & GeV/$c$ & vertex $|\vec q|$\\
34 & \code{nui} & GeV & vertex $\nu$\\
35 & \code{Q2i} & (GeV/$c$)$^2$ & vertex $Q^2$\\
36 & \code{Wi} & GeV & vertex $W$, computed with the struck-nucleon mass $M_N$ (no Fermi motion)\\
37 & \code{xbji} & -- & vertex $x$\\
38 & \code{epsiloni} & -- & vertex $\varepsilon$\\
39 & \code{Emi} & GeV & vertex $E_\mathrm{miss}=\nu+M_A-E_h$\\
40 & \code{Pmi} & GeV/$c$ & vertex $P_\mathrm{miss}=|\vec p_h-\vec q|$\\
41 & \code{missmass} & GeV & reconstructed missing mass on the struck nucleon at rest (\secr{sec:recon})\\
42 & \code{mmnuc} & GeV & reconstructed missing mass on the target nucleus (\secr{sec:recon})\\
43 & \code{phad} & GeV/$c$ & $P_h^\mathrm{rec}$\\
44 & \code{t} & GeV$^2$ & reconstructed $-t$ (\secr{sec:recon})\\
45 & \code{ti} & GeV$^2$ & vertex $-t$\\
46 & \code{z} & -- & $E_h^\mathrm{rec}/\nu^\mathrm{rec}$\\
47 & \code{zi} & -- & $E_h/\nu$ at the vertex\\
48 & \code{pt2} & GeV$^2$ & $(P_h^\mathrm{rec})^2\sin^2\theta_{pq}^\mathrm{rec}$\\
49 & \code{pt2i} & GeV$^2$ & $P_h^2\sin^2\theta_{pq}$ at the vertex\\
50 & \code{thetapq} & rad & $\theta_{pq}^\mathrm{rec}$\\
51 & \code{phipq} & rad & $\varphi_{pq}^\mathrm{rec}$\\
52 & \code{thetapqi} & rad & vertex $\theta_{pq}$\\
53 & \code{phipqi} & rad & vertex $\varphi_{pq}$\\
54 & \code{pmpar} & GeV/$c$ & $P^\mathrm{rec}_{m,\parallel}$\\
55 & \code{pmper} & GeV/$c$ & $P^\mathrm{rec}_{m,\perp}$\\
56 & \code{pmoop} & GeV/$c$ & $P^\mathrm{rec}_{m,\mathrm{oop}}$\\
57 & \code{fry} & cm & minus the vertical raster coordinate\\
58 & \code{radphot} & GeV & $E_{\gamma1}+E_{\gamma2}+E_{\gamma3}$\\
59 & \code{pfermi} & GeV/$c$ & $p_F$ with the sign of $\hat p_F\cdot\hat q$ (0 without Fermi motion)\\
60 & \code{siglab} & $\mu$b/(MeV sr$^2$) & $\sigma$ at the vertex, with the Coulomb factor (\secr{sec:xsec})\\
61 & \code{sigcm} & $\mu$b/(MeV$^2$ rad) & the centre-of-mass value written by the model (\secr{sec:xsec})\\
62 & \code{Weight} & see \secr{sec:xsec} & the event weight $w$\\
63 & \code{decdist} & cm & path flown by the hadron before its decay, or the whole path when it did not decay (\secr{sec:rad})\\
64 & \code{Mhadron} & MeV & mass of the particle that reached the detectors (the muon mass after a pion decay)\\
\end{simctable}

For pion production \code{siglab} is $\sigma$ of \eqr{eq:peepi} after the factors of \eqr{eq:cgfactors} and the Coulomb
factor, \code{sigcm} the blend of \eqr{eq:blend}, and \code{Weight} $w$ of \eqr{eq:weight} in $\mu$b/(MeV sr$^2$).
With \code{doing\_pizero} twelve columns follow (\tabr{tab:pi0columns}); with \code{using\_tgt\_field} eight more
(the angles to the target polarisation), after them.

\begin{simctable}{footnotesize}{L{0.06\linewidth}L{0.17\linewidth}L{0.11\linewidth}L{0.66\linewidth}}
\caption{The $\pi^0$ columns (one column per row).}\lbl{tab:pi0columns}\\
\toprule\rowcolor{simchead}\simcth{\#} & \simcth{name} & \simcth{unit} & \simcth{content}\\\midrule\endfirsthead
\toprule\rowcolor{simchead}\simcth{\#} & \simcth{name} & \simcth{unit} & \simcth{content}\\\midrule\endhead
\bottomrule\endlastfoot
65 & \code{xcal\_gamma1} & cm & photon 1 at the calorimeter, $x$ ($-10^{10}$ if missed)\\
66 & \code{ycal\_gamma1} & cm & photon 1 at the calorimeter, $y$ ($-10^{10}$ if missed)\\
67 & \code{Egamma1} & MeV & $E_1$ of \eqr{eq:pi0decay}\\
68 & \code{Pgamma1x} & MeV & $k_{1,x}$ (laboratory)\\
69 & \code{Pgamma1y} & MeV & $k_{1,y}$\\
70 & \code{Pgamma1z} & MeV & $k_{1,z}$\\
71 & \code{xcal\_gamma2} & cm & photon 2 at the calorimeter, $x$\\
72 & \code{ycal\_gamma2} & cm & photon 2 at the calorimeter, $y$\\
73 & \code{Egamma2} & MeV & $E_2$\\
74 & \code{Pgamma2x} & MeV & $k_{2,x}$\\
75 & \code{Pgamma2y} & MeV & $k_{2,y}$\\
76 & \code{Pgamma2z} & MeV & $k_{2,z}$\\
\end{simctable}

% =====================================================================================================================
\section{Validation in SIMC++}\lbl{sec:validation}
\begin{simctable}{footnotesize}{L{0.46\linewidth}L{0.2\linewidth}L{0.34\linewidth}}
\caption{CTest tests that cover pion production (one test per row).}\lbl{tab:tests}\\
\toprule\rowcolor{simchead}\simcth{test} & \simcth{labels} & \simcth{what it checks}\\\midrule\endfirsthead
\toprule\rowcolor{simchead}\simcth{test} & \simcth{labels} & \simcth{what it checks}\\\midrule\endhead
\bottomrule\endlastfoot
\code{l1\_peepi} & l1 & \eqr{eq:peepi} against the Fortran table\\
\code{l1\_sigmaid} & l1 & \eqr{eq:maid}\\
\code{l1\_sig\_multipole} & l1 & the unused multipole model\\
\code{l1\_transform\_to\_cm\_pion} & l1 & \eqsr{eq:boost}{eq:jold}\\
\code{l1\_pizero\_decay\_fixedrng} & l1 & \eqr{eq:pi0decay}\\
\code{l1\_pion\_data\_files} & l1 & the MAID and $\Delta$ tables (checksums)\\
\code{l1\_pion\_plugin} & l1 & the plug-in \code{pion} against \code{peepi}\\
\code{l1\_pion\_state} & l1 & the set-up of the pion runs (masses, limits, tables)\\
\code{l1\_generate\_pion\_fixedrng} & l1 & \secr{sec:gen} with a fixed random stream\\
\code{l1\_complete\_ev\_pion} & l1 & \secr{sec:complete}\\
\code{l1\_complete\_recon\_ev\_pion} & l1 & \eqr{eq:recmm}\\
\code{l1\_complete\_main\_pion} & l1 & \eqsr{eq:cgfactors}{eq:weight}\\
\code{l1\_replay\_event\_capture\_pion\_pip\_h} & l1, replay & every routine of captured H($\pi^+$) events\\
\code{l1\_simcrun\_replay\_pion\_pip\_h} & l1, replay & the complete run on a captured random stream, bitwise\\
\code{l1\_simcrun\_replay\_pion\_pim\_d} & l1 & the same, D($\pi^-$)\\
\code{l1\_montecarlo\_replay\_pion\_pip\_h} & l1, replay & the spectrometer calls of the capture\\
\code{l1\_project\_decay} & l1 & \eqsr{eq:decaypoint}{eq:decayproducts}\\
\code{l1\_transp\_decay\_hms} & l1 & decay inside the HMS magnets\\
\code{l1\_mc\_hms\_decay} & l1 & the HMS model with decay\\
\code{l1\_decay\_fixedrng} & l1 & the decay kinematics with a fixed random stream\\
\code{l1\_simcrun\_ranlux\_up\_pion\_h\_k1\_rad\_dec0\_A} & l1, replay, slow, upstream\_layout & full run, RANLUX, compat, against the upstream Fortran run\\
\code{l1\_simcrun\_ranlux\_up\_upf1\_pion\_c12\_wp0\_A} & l1, replay, upstream\_layout & the same, C-12 quasi-free $\pi^+$\\
\code{l1\_simcrun\_ranlux\_up\_upf1\_pion\_h\_wp3\_A} & l1, replay, slow, upstream\_layout & the same, H($\pi^-\Delta^{++}$)\\
\code{l1\_pi0calo\_state} & l1 & the $\pi^0$/calorimeter set-up\\
\code{calo\_selftest} & l1, unit & the calorimeter and the $\pi^0$ photons\\
\code{regression\_identity\_pion\_h\_k1\_rad\_dec1} & regression & compat-mode outputs against the baseline\\
\code{regression\_corrected\_pion\_h\_k1\_rad\_dec1} & regression & default-mode outputs against the baseline\\
\code{regression\_identity\_calo\_excl\_h\_ng2} & regression & $\pi^0$ with two photons, compat mode\\
\code{regression\_upstream\_u59\_p15f1\_pi0\_shms\_A} & regression\_upstream, slow & $\pi^0$ in the SHMS, against the upstream run\\
\code{regression\_shipped\_shipped\_test\_pion\_he\_pct} & regression\_shipped, slow & a shipped $A\ge3$ pion input\\
\code{fortran\_diff\_D038\_F10\_sigcm2\_stale} & fortran\_diff, unit & a default-mode correction of the model (\secr{sec:origin})\\
\end{simctable}
Recorded in the validation ledger (\code{project\_control/VALIDATION\_LEDGER.md}): the L1 tables of the model
(about 49\,000 rows, all columns bitwise, plug-in 5047 of 5047 rows; entry of 2026-10-06, task P9.1); the event tables,
both capture replays and seven L2 comparisons with the Fortran in the compat mode, all passing, the default mode
differing only by the central cross section and the hadron $\delta$ cut (entry of 2026-10-07, task P9.2). The
\code{fortran\_diff\_*} tests of the findings of \secr{sec:origin} are listed in FORTRAN\_DIFFERENCES.md.

% =====================================================================================================================
\section{Fortran origin and Fortran findings}\lbl{sec:origin}
\begin{simctable}{footnotesize}{L{0.42\linewidth}L{0.58\linewidth}}
\caption{Fortran routine $\to$ C++ function (one routine per row; the common routines are in \ctabr{tab:map}).}\lbl{tab:map}\\
\toprule\rowcolor{simchead}\simcth{Fortran (upstream 59ccdf6)} & \simcth{C++ (src/...)}\\\midrule\endfirsthead
\toprule\rowcolor{simchead}\simcth{Fortran (upstream 59ccdf6)} & \simcth{C++ (src/...)}\\\midrule\endhead
\bottomrule\endlastfoot
\file{dbase.f:138}-161 (pion flags, stops) & \csrc{core/RunSetup.cpp: DeriveReactionFlags}; \csrc{setup/SetupBuilder.cpp: BuildRunSetup}\\
\file{dbase.f:333}-393 (pion masses) & \csrc{setup/SetupReaction.cpp: BuildReactionAndKinematics}\\
\file{dbase.f:567}-591 (momentum distribution) & \csrc{setup/SetupTheory.cpp: ReadMomentumDistribution}\\
\file{dbase.f:598}-628 (spectral function) & \csrc{physics/nuclear/SpectralFunctionTables.cpp: SfTable::load}\\
\file{init.f:344}-513 (limits) & \csrc{setup/SetupLimits.cpp: BuildLimits}\\
\file{event.f} generate, l.276-391 & \csrc{reactions/pion/PionProduction.cpp: PionProduction::generate}\\
\file{sf\_lookup.f} generate\_em & \csrc{physics/nuclear/SpectralFunctionTables.cpp: SfTable::generateEm}\\
\file{event.f} complete\_ev, l.478-495 (table) & \csrc{reactions/pion/PionProduction.cpp: DeltaMassTable::load}\\
\file{event.f} complete\_ev, l.620-631 ($M_\Delta$) & \csrc{reactions/pion/PionProduction.cpp: DeltaMassTable::sample}\\
\file{event.f} complete\_ev, l.608-1054 & \csrc{reactions/pion/PionProduction.cpp: PionProduction::completeEv}\\
\file{event.f} complete\_ev, l.723-771 & \csrc{reactions/pion/PionProduction.cpp: PionProduction::pionMainKinematics}\\
\file{event.f} complete\_ev, l.796-897 & \csrc{reactions/pion/PionProduction.cpp: PionProduction::polarizedTargetAngles}\\
\file{pizero\_decay.f} pizero\_decay & \csrc{physics/pion/PizeroDecay.cpp: pizeroDecay}\\
\file{event.f} complete\_recon\_ev, l.1325-1354 & \csrc{reactions/EepReactionBase.cpp: EepReactionBase::completeRecon}\\
\file{event.f} complete\_main, l.1461-1503, 1561-1562 & \csrc{reactions/pion/PionProduction.cpp: PionProduction::completeMain}\\
\file{simc.f} calculate\_central, l.1245-1296 & \csrc{reactions/pion/PionProduction.cpp: PionProduction::prepareCentral}\\
\file{jacobians.f} transform\_to\_cm & \csrc{reactions/EventChain.cpp: transformToCm}\\
\file{physics\_pion.f} peepi & \csrc{physics/pion/PionCrossSections.cpp: peepi}; plug-in \csrc{physics/pion/PionModel.cpp: PionModel::sigccMain}\\
\file{physics\_pion.f} sig\_param\_2021 & \csrc{physics/pion/PionCrossSections.cpp: sigParam2021}\\
\file{physics\_pion.f} exclfit & \csrc{physics/pion/PionCrossSections.cpp: exclfit}\\
\file{physics\_pion.f} sigmaid (tables) & \csrc{physics/pion/PionCrossSections.cpp: MaidTables::load}\\
\file{physics\_pion.f} sigmaid & \csrc{physics/pion/PionCrossSections.cpp: sigmaid}\\
\file{physics\_pion.f} sig\_blok, sig\_param04, sig\_param\_3000 & \csrc{physics/pion/PionCrossSections.cpp: sigBlok, sigParam04, sigParam3000}\\
\file{physics\_pion.f} sig\_multipole & \csrc{physics/pion/SigMultipole.cpp: SigMultipole}\\
\file{radc.f:297}-308 (tail 1, case 4) & \csrc{radiation/Radc.cpp: generateRad}\\
\file{simc.f:1500}-1566 ($\pi^0$ photons) & \csrc{driver/PizeroCalorimeter.cpp: PizeroCalorimeter::detect}\\
\file{shared/project.f} project & \csrc{optics/CosyEngine.cpp: project}\\
\file{shared/transp.f} transp (decay) & \csrc{optics/CosyEngine.cpp: CosyEngine::transp}\\
\file{results\_write.f:178}-286 & \csrc{output/NtupleLayout.cpp: FillUpstreamNtuple}\\
\file{simc.f:394}-397 (genvol) & \csrc{reactions/EepReactionBase.cpp: EepReactionBase::genVolume}\\
\end{simctable}

\begin{simcfindings}{footnotesize}{L{0.18\linewidth}L{0.5\linewidth}L{0.32\linewidth}}
\caption{Fortran findings that touch pion production (one finding per row; the common ones are in \ctabr{tab:findings}).}\lbl{tab:findings}\\
\toprule\rowcolor{simchead}\simcth{id} & \simcth{Fortran behaviour} & \simcth{default / compat}\\\midrule\endfirsthead
\toprule\rowcolor{simchead}\simcth{id} & \simcth{Fortran behaviour} & \simcth{default / compat}\\\midrule\endhead
\bottomrule\endlastfoot
\finding{U-A3} & 59ccdf6: \code{which\_pion} 2/3 on $A\ge3$ weighted per nucleon, without $Z$ (weight $\times1/Z$ against 294f662) & 294f662 weight in both modes (author asked)\\
\finding{D-038 F1} & H($\pi^-\Delta^{++}$) stopped (fixed upstream in 20d46fc) & runs / runs (294f662 target: stop)\\
\finding{D-038 F2} & $\pi^0$ from the neutron accepted on hydrogen & refused / accepted\\
\finding{D-038 F5} & central cross section with $\varepsilon=-t=\varphi_{pq}=0$ & central kinematics / Fortran\\
\finding{D-038 F10} & MAID value of an earlier event kept for $W\ge2$ GeV (not a tree column) & 0 / Fortran\\
\finding{D-038 F15} & \texttt{.hist} final-state line missing for \code{which\_pion} 11 & printed / Fortran\\
\finding{D-038 F3} & coherent production: the input final-nucleus mass is never used & kept in both\\
\finding{D-038 F4} & coherent production with the nucleon model, $n_\mathrm{tgt}=Z$ or $N$, negative cross sections possible & kept in both\\
\finding{D-038 F9} & a tail-1 photon draws a second $\Delta$ mass; the \texttt{.hist} shows the last one & kept in both\\
\finding{D-038 F12} & unused multipole model in single precision with stale amplitudes & reproduced\\
\finding{D-038 F11} & relative momentum of the recoil pair: stale local value & not an output\\
\finding{D-038 F13} & unused fit with a pole at $|t|=0.02$ GeV$^2$ & dead code\\
\finding{D-038 F14} & MAID: nearest lower bin, no interpolation, 6 of 23 angle rows & the coded model\\
\finding{D-038 pi0 decay} & decay failure branch unreachable & reproduced\\
\finding{D-038 F16} & every STOP exits with status 0 & exit status 2 or 3\\
\finding{D-062 G1} & $\pi^0$ photons decayed with the $\pi^0$ mass: STOP with \code{doing\_decay} & massless, stable / Fortran\\
\finding{D-062 G3} & second photon starts with the first photon's path length & each from 0 / Fortran\\
\finding{D-062 G4} & focal-plane columns of the last photon that hit; no $\pi^0$ reconstruction & reproduced\\
\finding{D-062 G8} & a tail-1 photon draws a second photon pair & kept in both\\
\finding{D-065 P1} & $\pi^0$ transported through a magnetic spectrometer & kept in both (no decay in default)\\
\finding{D-067 T13} & $\pi^0$ bent in the target field & kept in both\\
\finding{D-067 T14} & field + $\pi^0$: photon columns overwrite the field columns (294f662 layout) & own columns / own columns\\
\finding{U-A4} & field + $\pi^0$: 84 values in arrays of 80 (59ccdf6) & full layout in both\\
\finding{D-037 F7} & remaining drift after a decay computed with the muon slopes & pion slopes / Fortran\\
\finding{D-037 F8} & decay inside a magnet applied at its entrance or exit & kept in both\\
\finding{D-037 F9} & decay of a non-$\pi$/K hadron stops the run & refused at input / stop\\
\finding{D-006 pass\_cuts} & hadron $\delta$ upper cut with the electron limit & hadron limit / Fortran\\
\end{simcfindings}

\begin{simcfindingbox}
\textbf{Root of the squared energy equation.} \eqr{eq:quad} squares $\tau-2cE_h=2aP_h$ and the code does not check
the sign afterwards. A numerical check of the formula alone (Python, H($\pi^+$)n, $4\times10^5$ random $(\nu, Q^2,
\theta_{pq})$ points, $\nu<8$ GeV) found that every accepted root with $\theta_{pq}>90^\circ$ violated momentum
conservation (recoil mass $\ne M_X$), and none with $\theta_{pq}<90^\circ$. Usual set-ups keep $\theta_{pq}$ small,
so runs are probably not affected; not checked on reference runs \fsrc{event.f:676-698}.
\end{simcfindingbox}

\begin{simcfindingbox}
\textbf{Generation limits with the nominal $\Delta$ mass.} The electron-energy and photon limits of
\eqsr{eq:vertexedge}{eq:sumegen} are computed once with $M_X=M_\Delta=1232$ MeV \fsrc{dbase.f:361-367; init.f:399-462},
while the event takes $M_X$ from 1028 to 1627 MeV (58\,\% below 1232 MeV, \eqr{eq:deltamass}). When the
energy-conservation term is the active upper limit of $\Sigma_\mathrm{max}$, events with a light $\Delta$ and the
largest electron energies are never generated. Found by reading the code; the size of the effect was not measured.
\end{simcfindingbox}

\begin{simcfindingbox}
\textbf{Coherent $\pi^0$.} \code{doing\_pizero} with \code{which\_pion} 10 or 11 passes the set-up checks for heavy
targets, but the $\pi^0$ branch of the mass set-up has no case for them \fsrc{dbase.f:334-351}: $M_N$ stays 0 and $M_X$
is set to the $A-1$ guess $M_A-M_p$. SIMC++ reproduces this (the model then refuses $M_N=0$ at start-up). Found by
reading the code; not run in the Fortran.
\end{simcfindingbox}

\begin{simcrefs}
\simcref{maid} D. Drechsel, S. S. Kamalov, L. Tiator, Eur.\ Phys.\ J.\ A \textbf{34}, 69 (2007): MAID2007 (the code says
  ``MAID 2007'' \fsrc{physics\_pion.f:580}, without a reference).
\simcref{brauel} P. Brauel et al., Z.\ Phys.\ C \textbf{3}, 101 (1979) (cited at \fsrc{physics\_pion.f:408}).
\simcref{cgln} G. F. Chew, M. L. Goldberger, F. E. Low, Y. Nambu, Phys.\ Rev.\ \textbf{106}, 1345 (1957) (cited at
  \fsrc{physics\_pion.f:333}).
\simcref{walker} R. L. Walker, Phys.\ Rev.\ \textbf{182}, 1729 (1969) (cited at \fsrc{physics\_pion.f:343}).
\simcref{bartl} E. Bartl, Nucl.\ Phys.\ B \textbf{62}, 267 (1973) (cited at \fsrc{physics\_pion.f:351}).
\simcref{devenish} R. C. E. Devenish, D. H. Lyth, Phys.\ Rev.\ D \textbf{5}, 47 (1972) (cited as ``Phys. Rev. C D5'' at
  \fsrc{jacobians.f:189} for the old jacobian).
\end{simcrefs}
\end{document}
